\Subhead{2.6.1}{Электромагнитные волны в металлических волноводах}

Пусть идеальный проводящий волновод однороден вдоль $z$, а поля зависят как $\e^{\ii\beta z-\ii\omega t}$. Продольные компоненты удовлетворяют поперечному уравнению Гельмгольца:
\begin{equation}
 (\Delta_\perp+k_\perp^2)\{E_z,H_z\}=0,\quad
 \beta^2=k^2-k_\perp^2,\quad k=\omega\sqrt{\varepsilon\mu}.
\end{equation}
Для TM-мод ($H_z=0$) $E_z$ обращается в нуль на металлической стенке. Для TE-мод ($E_z=0$) нормальная производная $H_z$ равна нулю. TEM-мода имеет обе продольные компоненты нулевыми; в односвязном полом металлическом волноводе без внутренних проводников она невозможна, поскольку поперечная электростатическая задача с одним постоянным потенциалом даёт нулевое поле.

В прямоугольнике $a\times b$
\begin{equation}
 k_{\perp,mn}^2=(m\pi/a)^2+(n\pi/b)^2,\quad
 \omega_c=\frac{k_\perp}{\sqrt{\varepsilon\mu}}.
\end{equation}
У TM оба индекса ненулевые, у TE один может быть нулём, но не оба. При $a>b$ основная мода TE$_{10}$. Ниже отсечки $\beta$ мнима и мода не переносится как бегущая волна в бесконечном однородном участке. Выше отсечки
\begin{equation}
 v_{\rm ph}=\omega/\beta>v,\qquad v_g=\frac{v^2\beta}{\omega}<v,
 \quad v=1/\sqrt{\varepsilon\mu},
\end{equation}
если материал недиспергирующий. Сверхсветовая фазовая скорость не означает сверхсветовой передачи сигнала. Модальные отношения поперечных полей различны: $Z_{\rm TE}=\omega\mu/\beta$, $Z_{\rm TM}=\beta/(\omega\varepsilon)$.
